% -*- mode: fundamental -*-

% Formal spec of BSV semantics
% Started by Nikhil, April 23, 2015

\documentclass[11pt]{article}

% ================================================================
\usepackage[latin1]{inputenc}
\usepackage[T1]{fontenc}
\usepackage{latexsym}
\usepackage{makeidx}
\usepackage{alltt}
\usepackage{verbatim}
\usepackage{fancyvrb}
% \usepackage{moreverb}
\usepackage{ae}
\usepackage{aecompl}

  \usepackage[pdftex,colorlinks=true,bookmarksopen, pdfstartview=FitH,
              linkcolor=blue, citecolor=blue, urlcolor=blue]{hyperref}
  \pdfcompresslevel=9
  \usepackage[pdftex]{graphicx}

% ================================================================

% HORIZONTAL MARGINS
% Left margin, odd pages: 1 inch ("0" + 1)
\setlength{\oddsidemargin}{0in}
% Left margin, even pages: 1 inch ("0" + 1)
\setlength{\evensidemargin}{0in}
% Text width 6.5 inch (so other margin is 1 inch).
\setlength{\textwidth}{6.5in}
% ----------------
% VERTICAL MARGINS
% Top margin 0.5 inch (-0.5 + 1)
\setlength{\topmargin}{-0.5in}
% Head height 0.25 inch (where page headers go)
\setlength{\headheight}{0.25in}
% Head separation 0.25 inch (between header and top line of text)
\setlength{\headsep}{0.25in}
% Text height 9 inch (so bottom margin 1 in)
\setlength{\textheight}{9in}
% ----------------
% PARAGRAPH INDENTATION
\setlength{\parindent}{0in}
% SPACE BETWEEN PARAGRAPHS
\setlength{\parskip}{\medskipamount}
% ----------------
% STRUTS
% HORIZONTAL STRUT.  One argument (width).
\newcommand{\hstrut}[1]{\hspace*{#1}}
% VERTICAL STRUT. Two arguments (offset from baseline, height).
\newcommand{\vstrut}[2]{\rule[#1]{0in}{#2}}
% ----------------
% HORIZONTAL LINE ACROSS PAGE:
\newcommand{\hdivider}{\noindent\mbox{}\hrulefill\mbox{}} 
% ----------------
% EMPTY BOXES OF VARIOUS WIDTHS, FOR INDENTATION
\newcommand{\hm}{\hspace*{1em}}
\newcommand{\hmm}{\hspace*{2em}}
\newcommand{\hmmm}{\hspace*{3em}}
\newcommand{\hmmmm}{\hspace*{4em}}
% ----------------
% VARIOUS CONVENIENT WIDTHS RELATIVE TO THE TEXT WIDTH, FOR BOXES.
\newlength{\hlessmm}
\setlength{\hlessmm}{\textwidth}
\addtolength{\hlessmm}{-2em}

\newlength{\hlessmmmm}
\setlength{\hlessmmmm}{\textwidth}
\addtolength{\hlessmmmm}{-4em}
% ----------------
% ``TIGHTLIST'' ENVIRONMENT (no para space betwee items, small indent)
\newenvironment{tightlist}%
{\begin{list}{$\bullet$}{%
    \setlength{\topsep}{0in}
    \setlength{\partopsep}{0in}
    \setlength{\itemsep}{0in}
    \setlength{\parsep}{0in}
    \setlength{\leftmargin}{1.5em}
    \setlength{\rightmargin}{0in}
    \setlength{\itemindent}{0in}
}
}%
{\end{list}
}
% ----------------
% ITALICISE WORDS
\newcommand{\ie}{\emph{i.e.,}}
\newcommand{\eg}{\emph{e.g.,}}
\newcommand{\Eg}{\emph{E.g.,}}
\newcommand{\etc}{\emph{etc.}}
\newcommand{\via}{\emph{via}}
\newcommand{\vs}{\emph{vs.}}

\newcommand{\SV}{SV}
\newcommand{\synscemidpi}{Synthesizable SCE-MI DPI}

% ----------------
% CODE FONT (e.g. {\cf x := 0}).
\newcommand{\cf}{\footnotesize\tt}

% ----------------------------------------------------------------
% BNF grammar

\newcommand{\te}[1]{\texttt{#1}}
\newcommand{\nterm}[1]{\emph{#1}}
\newcommand{\term}[1]{\texttt{#1}}
\newcommand{\many}[2]{#1 #2 ... #2 #1}
\newcommand{\opt}[1]{[ #1 ]}
\newcommand{\alt}{{$\mid$}}
\newcommand{\gram}[2]{    \hm\makebox[10em][l]{\it #1}\makebox[1.5em][l]{::=}    #2}
\newcommand{\grammore}[1]{\hm\makebox[10em][l]{      }\makebox[1.5em][l]{}       #1}
\newcommand{\gramalt}[1]{ \hm\makebox[10em][l]{      }\makebox[1.5em][l]{\alt}   #1}

% ----------------------------------------------------------------
% HERE BEGINS THE DOCUMENT

% ================================================================

\begin{document}

% ================================================================
% Title etc.

\begin{center}

\centerline{\includegraphics[width=2.5in,angle=0]{Figures/Fig_Logo_2color_on_white}}

\vspace{2ex}

{\Large\bf Formal Specification of BSV's Elaboration and Dynamic Semantics}

\vspace{1ex}

Rishiyur S. Nikhil

\vspace{1ex}

Bluespec, Inc. \\
July 28, 2015 (rev. July 15, 2016)

\end{center}

% ================================================================

\begin{center}
\begin{minipage}{\hlessmm}
\begin{center}
\bf Abstract
\end{center}

BSV is a High Level Hardware Design Language (HLHDL), intended for the
full spectrum of digital hardware designs, from signal-processing
accelerators to CPUs, SoCs and all kinds of Intellectual Property
Blocks (IP Blocks).  BSV has been in use commercially and in academia
since the early 2000s.

\hm Although BSV was inspired by formally-defined languages
(specifically the Haskell functional programming language and Term
Rewriting Systems), its formal elaboration and dynamic semantics have
not been properly documented to date.  So far, this has not been an
issue, but there is now a growing interest in using BSV for formally
verified hardware, because of its clean and high level semantics.

\hm This document provides a formal specification of BSV elaboration
and dynamic (execution) semantics. The specification is written as a
Haskell program and is therefore executable.  This document is
accompanied by the full Haskell code as well as the (kernel) BSV
source code for the fully-worked out examples which are described in
this document.

\end{minipage}
\end{center}

% ================================================================

\section{Introduction and History}

BSV~\cite{Bluespec2015} is a High Level Hardware Design Language
(HLHDL), intended for the full spectrum of digital hardware designs,
from signal-processing accelerators to CPUs, SoCs and all kinds of
Intellectual Property Blocks (IP Blocks).  BSV has been in use
commercially and in academia since the early 2000s.

\begin{tabbing}
\hmm \= \hspace*{1.5in} \= \kill
     \> \emph{Acronyms} \\
     \> HDL   \> Hardware Description Language \\
     \> HLHDL \> High Level Hardware Description Language \\
     \> HLS   \> High Level Synthesis \\
     \> HW    \> Hardware \\
     \> PL    \> Programming Language \\
     \> SoC   \> System on a Chip \\
     \> SW    \> Software \\
     \> TRS   \> Term Rewriting Systems
\end{tabbing}

% ----------------------------------------------------------------

\subsection{Historical context}

From the beginning, BSV has been inspired by formally-defined
languages, specifically Term Rewriting Systems~\cite{Klop1992a,
Terese2003, Baader98a} (TRSs) and the Haskell functional programming
language~\cite{PeytonJones2003a}.  It started with research by Shen
and Arvind at MIT to use Term Rewriting Systems to formally specify
and verify complex concurrent hardware, such as deep processor
pipelines and distributed coherent caches for multi-core
systems~\cite{Arvind1998a,StoyFME2001}.  Similar approaches (using
rewrite rules to specify complex concurrency) may also be found in
TLA+~\cite{Lamport2002a} and its use to verify cache coherence
protocols~\cite{Akhiani1999a,Joshi2003a}, in UNITY~\cite{Chandy1988a}
and in Event-B~\cite{Metayer2005a}.

From the specification research of Shen and Arvind emerged the idea of
directly \emph{synthesizing} hardware from TRS specifications.  This
idea was developed significantly by Hoe and Arvind~\cite{Hoe2000a,
Hoe2000b,Hoe2004a} at MIT.

In 2000, activity on productizing these ideas moved from MIT into
Sandburst Corporation.  There, Lennart Augustsson \emph{et al.}
designed and implemented a new language called
Bluespec~\cite{Augustsson2001a}.  Whereas the previous prototype from
MIT had only rudimentary circuit-description facilities, focusing more
on solving the problem of synthesis from TRSs, the new language
borrowed heavily from Haskell to also provide a powerful circuit
description language (also called ``static elaboration'' in HDL
jargon).  It had Haskell-like syntax, type definitions, polymorphism,
type classes, higher-order functions, total absence of side effects,
laziness, and a monadic structure for static elaboration.  Although it
was an industrial-strength language, it was only intended for internal
use inside Sandburst.  Nowadays we refer to this second-generation
language as ``Bluespec Classic''.

In 2003, activity moved again, to Bluespec, Inc., with an aim to make
it a publicly available Hardware Design Language (HDL).  In this third
generation, first released in 2004, its syntax was changed from
Haskell's very spare syntax to something more similar to
SystemVerilog~\cite{IEEESystemVerilog2012a} which is more familiar to
hardware design engineers, the primary audience. The name of the
language was changed to \emph{BSV} (Bluespec SystemVerilog).  Support
for multiple clock domains was added, with statically typed clocks and
resets and clock domain discipline.  But apart from this syntax change
and addition of clocks, the underlying language and its basic
semantics has remained the roughly same (even today you can
successfully compile and run BSV programs from 2004!).

% ----------------------------------------------------------------

\subsection{Why no earlier formalization?}

\noindent Like many industrial-strength languages, BSV is not a tiny
language (although the language-definition part of the Reference
Guide~\cite{Bluespec2015} is barely 150 pages, compared to more than
1300 pages for SystemVerilog~\cite{IEEESystemVerilog2012a}), and so
defining formal semantics in gory detail is not trivial (and may not
be readable by anyone!).  BSV's semantics is described informally in
the Reference Guide, in a number of books and papers (e.g.,
\cite{Nikhil2008a, Nikhil2010a, Nikhil2011b, Nikhil2013}), and in
Bluespec's BSV training materials.  By and large this has been
perfectly adequate for practical use.

There is now a strong emerging interest in \emph{mechanized} formal
verification of computing systems that include components implemented
in hardware.  The interest spans CPUs, memory systems and SoCs as well
as hardware accelerators, hardware implementations of cryptography,
and so on.  There are many groups---commercial, non-profit research
and academic---who are keen on using BSV for this purpose because it
is both industrial-strength and has a very clean and high level
semantics.  This activity requires formal specification of BSV, hence
this document.

% ----------------------------------------------------------------

\subsection{Document overview}

This document is accompanied separately by Haskell code that specifies
the semantics.  The Haskell code is executable, and comes with several
example ``kernel BSV'' programs on which it can be run.

That Haskell code, and not this document, should be regarded as the
official formal specification.  This document should be regarded as a
reading guide to that Haskell code.  The Haskell code fragments shown
here are from the accompanying Haskell code.  In case they are out of
sync, the accompanying Haskell code should be taken as definitive.

\emph{Note (July 27, 2015): This is the first public release of this
document and the Haskell code.  It is quite possible that both the
semantics and this document need fixes and/or refinement.  The author
gladly welcomes comments and suggestion from readers.}

A detailed table of contents follows next.  Then,
Sec.~\ref{sec_prelims} provides some general commentary about our
approach.  Sec.~\ref{sec_syntax} describes the concrete and abstract
syntax of our kernel language.  Sec.~\ref{sec_parsing} talks briefly
about parsing from concrete to abstract syntax.
Sec.~\ref{sec_static_elaboration} describes static elaboration of the
source program into a collection of module instances.

Sec.~\ref{sec_dynamic_semantics}, ``Dynamic Semantics'', is the heart
of this document.  Readers who are experienced at reading Haskell
and/or formal semantics of languages may wish to skip preceding
sections and go there directly.

Sec.~\ref{sec_examples} presents a set of four example kernel BSV
programs and their semantic interpretation (included with the
accompanying Haskell code).  The first two are explained in detail and
the latter two are left as exercises for the reader (although we
provide Haskell execution transcripts for all of them).
Sec.~\ref{sec_conclusion} concludes with some comparisons to other
hardware description languages and with a discussion of the
relationship of this formal semantics to the full BSV language and its
compilation with the \emph{bsc} compiler.

% ================================================================

\renewcommand*{\contentsname}{}

\section{Table Of Contents}

\tableofcontents

\hspace*{1em}

% ================================================================

\section{Preliminary comments}

\label{sec_prelims}

Certain aspects of BSV's formal semantics may seem novel to those
already familiar with semantics of programming languages or TRSs.  We
discuss them in the following sections.

% ----------------------------------------------------------------

\subsection{Static elaboration and static vs. dynamic semantics}

All SW programming languages have a dynamic control
structure.\footnote{The only exceptions are ancient languages like
early FORTRAN and perhaps a few languages for embedded programs.}
Specifically, \emph{frames} are pushed and popped on a stack as we
dynamically call and return from functions, methods and exceptions.

In hardware, there is no such dynamic control structure.  Hardware
consists of physical electronics; they do not grow or shrink during
execution!\footnote{Of course it is always possible to
\emph{represent} a dynamic stack in memory in a hardware system, but
that is a level of abstraction and interpretation above our current
level of discourse, namely physical hardware.}

Nevertheless, you will find functions and procedures in all modern
HDLs.  In BSV you will also see extensive use of methods, loops, and
recursion.  These are merely linguistic devices for \emph{circuit
description}, i.e., they are unfolded, or inlined by the HDL
compiler/interpreter to describe a circuit.  It is almost like
statically pre-building the possible tree of stack frames for a SW
program just by looking at its text.

In SW programming languages, \emph{static semantics} usually refers to
syntax, scoping and other well-formedness checks, type-checking, and
perhaps some simple syntax desugaring.  The ``output'' of static
semantics is still a program text that may be recognizably similar to
the original program text.  Everything after that is described by
\emph{dynamic semantics}, including dynamic control behavior like
conditionals, loops, function calls and returns, recursion,
exceptions, and so on.

In HDLs, \emph{static semantics} covers the same topics as in SW PLs,
but goes further to encompass static elaboration, i.e., the unfolding
of linguistic constructs like function application, method calls,
elaboration loops and recursion, and so on.  The result of static
elaboration is a hardware circuit.  This hardware circuit may have a
logical hierarchical structure, i.e., a top-level module containing
instantiated sub-modules which, in turn, may contain next-level
instantiated sub-modules, and so on.  At the leaves of this hierarchy
will be primitive modules such as registers, FIFOs, RAMs, and I/O
devices.  All the modules are connected by \emph{combinational
circuits}.  In HDLs, \emph{dynamic semantics} concerns the behavior of
this statically elaborated circuit.

In this document, we will not discuss the part of BSV static semantics
that coincides directly with Haskell (specifically types,
polymorphism, type classes, type-checking and monadic elaboration).
We will describe BSV's static elaboration, because modules, module
instances, and module \emph{sharing} are fundamental concepts.

% ----------------------------------------------------------------

\subsection{TRSs, Linear Orders, and Clocks}

Classical TRSs are non-deterministic, and have no notion of clocks.
The basic dynamic semantics of a TRS is almost trivial:
\begin{itemize}

\item Given the current state (a ``term''), find a part of the state
(a subterm) that matches the left-hand side of one of the rewrite
rules.  (There may be many such matches; choose one non-deterministically.)

\item Replace that subterm by the right-hand side of the rule (with
suitable variable substitutions), producing the next state.  We call
this a \emph{rule firing}.

\item Repeat and rinse.

\end{itemize}

\noindent The description above is sequential (one rule at a time) and
may seem at odds with our earlier claim that this is a favorite way
amongst researchers to model complex concurrency.  The apparent
paradox is resolved by realizing (a) it is non-sequential in that
there is no prescribed ``control flow'' or ``program counter'' to
sequence rules; it is purely reactive--which rule can fire depends
only on the instantaneous state, and (b) \emph{any} actual concurrent
execution that preserves the logical sequentiality of the semantics is
acceptable (\emph{linearizability}).  Another way of saying this is
that each rewrite is an \emph{atomic transaction} on the
state~\cite{Lampson81}.

BSV goes further and refines TRS semantics to include \emph{clocks},
which are fundamental to most digital HW\footnote{Digital circuits can
also be \emph{asynchronous} or clock-free.  It is a fascinating
question whether a BSV-like language can be compiled to asynchronous
hardware, but we do not explore that further here.}:

\begin{itemize}

\item Every BSV rule fires ``within'' a clock.\footnote{This is a
design choice based on current tractability; a more general future
approach may allow rules to use multiple clocks.}

\item Many BSV rules may fire within a clock, but there will always be
a \emph{logical linear order} (serialization) that explains their behavior.

\item There will be practical constraints (described later) that
preclude certain rules from firing within the same clock in a certain
order.

\end{itemize}

% ----------------------------------------------------------------

\subsection{Haskell executability}

Our specification is given as a Haskell program, which is of course
executable.  It can be characterized as an abstract operational
semantics (high-level, implementation-independent).

% ----------------------------------------------------------------

\subsection{Kernel language only}

As is common in the literature, we will describe BSV's semantics for a
\emph{kernel} or subset of the language.  Extending it to the full
language would be routine but tedious to describe and to understand,
without providing any useful additional illumination.  Nevertheless,
our kernel language is more comprehensive than previous toy
descriptions of BSV semantics, and captures the essence of actual BSV
code behavior (specifically: module instances and module sharing).
Further, it is \emph{the} mental model that BSV designers use in
practice when developing and debugging BSV code.

% ----------------------------------------------------------------

\subsection{Rule firing constraints due to module sharing}

As mentioned earlier, we will be introducing constraints that prevent
certain rules from firing within a clock in a certain order.  Many of
these constraints have functional motivations.  For example, whether a
FIFO logically enqueues before it dequeues, or vice versa, can affect
its functionality in certain corner cases (when it is full or empty).

Other constraints arise due to \emph{module sharing}.  In digital
hardware, a wire can be driven with only one, stable value (``0'' or
``1'') during each clock.  Consider two independent rules that each
occasionally drive a particular input wire of a shared module.  Those
rules therefore cannot fire in the same clock.  Our semantics will
capture these constraints that arise due to module sharing.

% ----------------------------------------------------------------

\subsection{The semantics does not rely on any built-in primitives}

This semantics does not prescribe any particular predefined set of
primitives (such as registers).  For example, it is perfectly feasible
with this semantics to define a version of BSV based entirely on FIFOs
or RAMs, or even registers with different semantics.  This semantics
is, in effect, parameterized by a set of primitives.  All it assumes
is that each primitive is supplied with:

\begin{itemize}

\item An abstract function for the the primitive semantics of each
method when invoked (within an Action),

\item the parallel (intra-rule) constraints on its methods (i.e., for
each pair of methods, whether or not those two methods can be invoked
in the same Action),

\item the concurrent (inter-rule) ordering constraints on its methods
(i.e., for each pair of methods, the logical ordering of those methods
if they are invoked in different rules).
\end{itemize}

Of course in practice our primitives include registers, Concurrent
Registers, FIFOs, BRAMs and more.  The accompanying Haskell code
includes registers and Concurrent Registers (CRegs).

% ----------------------------------------------------------------

\subsection{Schedule-construction is an orthogonal question and is not addressed here}

\label{sec_schedules}

We use the term \emph{schedule} to mean a linear sequence of rule
instances.  In each clock, we will execute rule instances according to
such a schedule (execution of a rule instance may be a no-op, because
of conflicts or because the rule is not enabled).

It is important to distinguish two quite separate questions:

\begin{itemize}

\item How should one to construct a schedule from the rules in a BSV
program?  This could range from a completely fixed, static schedule
chosen by a compiler (e.g., current \emph{bsc}) to a completely
dynamic function that selects the next rule for execution just before
executing it.

\item Given a schedule, what is a \emph{correct execution} of the schedule?

\end{itemize}

In this document, we only address the latter question.  A scheduling
function is provided as an argument to the main semantic function.
Thus, the semantics defined here can be seen as \emph{universally
quantified over all possible schedules} i.e., describing correct
execution for any possible schedule.

In this semantics there is no such thing as a ``wrong'' schedule.  Of
course, from a practical point of view, some schedules are likely to
be better than others because, for example, they admit more
concurrency (more rules can fire in a clock).  This semantics takes no
position on how to construct such a schedule.

Our description here is therefore more general than the current
\emph{bsc} compiler for BSV, which:

\begin{tightlist}

\item Only considers each rule once during each clock.

\item Statically fixes a priority (linearization) of the rules.

\end{tightlist}

Both these constraints are motivated purely by implementation
efficiency concerns.  In contrast, the semantic description here
admits schedules that may vary from clock to clock, and may include
duplicate rule execution within a clock.

% ================================================================

\section{Kernel Language Concrete and Abstract Syntax}

\label{sec_syntax}

The kernel language syntax is described below both as abstract syntax
(a Haskell data type) and as concrete syntax.  The parser reads concrete syntax from
a text file and produces a Haskell data structure corresponding to the
abstract syntax.

In the accompanying Haskell code, the file {\tt AST.hs} defines the
abstract syntax data types.  The file {\tt Parser.hs} implements the
parser using Haskell's \term{Parsec} parser combinators.

Metasyntax: we use conventional BNF below for concrete syntax, with
the following conventional extensions:
\begin{tabbing}
\hmm \= \emph{x} ... \emph{x} \hm \= means zero or more repetitions of \emph{x} \kill
\hmm \> \{ \emph{x} \}            \> means zero or more repetitions of \emph{x} \\
\hmm \> \emph{x} ... \emph{x}     \> means zero or more repetitions of \emph{x} \\
\hmm \> [  \emph{x} ]             \> means zero or one repetitions of \emph{x} (i.e., \emph{x} is optional)
\end{tabbing}

\term{Typewriter} font is used for terminals, and \nterm{italics} for non-terminals.

% ----------------

\subsection{Programs}

A kernel BSV program is a list of module definitions followed by a
schedule.

{\bf Concrete syntax:}

\gram{Program}{\many{\nterm{ModuleDefinition}}{\hm} \nterm{Schedule}}

{\bf Abstract syntax:}

\begin{Verbatim}[frame=single, commandchars=\\\{\}]
type Program  = ([Binding], Schedule)
\end{Verbatim}

In the abstract syntax, each module definition is just a binding of
the module name to a lambda-expression whose arguments are the module
parameters and whose body is a module expression.

There must be exactly one module definition called \term{main}, and it
must have zero parameters.  Static elaboration is performed by
applying this niladic lambda expression (to an empty argument list).

Static elaboration will recursively instantiate modules; this is
sometimes called the module hierarchy.  Thus, each module instance has
a unique hierarchical name, which is a sequence of identifiers that
identify a path down this recursion, starting with \term{main}.

The hierarchical name of any module instance's component (such as a
rule or method instance) is simply the hierarchical name of its
enclosing module instance appended with component's name.

\begin{Verbatim}[frame=single, commandchars=\\\{\}]
type Ide      = String      -- Identifiers
type RuleName = String      -- Rule names
type MethName = String      -- Method names
type HierName = [Ide]       -- Hierarchical name
\end{Verbatim}

The schedule is a list of hierarchical names, each referring to a
unique post-elaboration instance of a rule.

\begin{Verbatim}[frame=single, commandchars=\\\{\}]
type Schedule   = [HierName]  -- list of rule instances
\end{Verbatim}

Having a schedule at the end of the source file is just a convenience
that allows us easily to exercise the semantics.  As described in
Sec.~\ref{sec_schedules}, this semantics takes no position on schedule
construction, just on correct execution of any given schedule.  One
could write variants of this Haskell code where the schedule is
constructed from analysis of the module definitions, or uses an run
time oracle, etc.  Some of our examples contain multiple schedules at
the end of the source file, and we comment-out all but one of them for
each execution.

% ----------------

\subsection{Bindings}

A binding associates an identifier with an expression/value.

{\bf Concrete syntax:}

\gram{Binding}{\term{let}  \nterm{Ide} \term{=}  \nterm{Expr}}

{\bf Abstract syntax:}

\begin{Verbatim}[frame=single, commandchars=\\\{\}]
type Binding = (Expr, Expr)
\end{Verbatim}

The lhs uses \term{Expr} instead of \term{Ide} to allow for possible
future extension where it could be an array selection, struct
selection or pattern.

% ----------------

\subsection{Module Definitions}

\label{sec_syntax_module_definitions}

The word ``module'' is used in quite different ways in different
languages and often concerns mechanisms for separate compilation and
namespace management, without any serious role in dynamic semantics.
In BSV, as in many HDLs, a module is a fundamental dynamic semantic
entity, representing a hardware unit with an interface.  Source
programs contain module expressions which evaluate to modules.
Elaborated programs contain module instances.  A single module
expression may result in multiple module instances.  A module instance
may be \emph{shared}, which can affect dynamic semantics.

A module expression is an expression which contains a sequence of
bindings, rules, and interface method definitions.

{\bf Concrete syntax:}

\gram{ModuleDefinition}{\term{module} \nterm{Ide} \opt{\term{\# (} \many{\nterm{Ide}}{\term{,}} \term{)}} \term{;}} \\
\grammore{\hmm \many{\nterm{Binding}}{\term{;}}} \\
\grammore{\hm \term{rules}} \\
\grammore{\hmm \many{\nterm{Rule}}{\hm}} \\
\grammore{\hm \term{methods}} \\
\grammore{\hmm \many{\nterm{MethodDef}}{\hm}} \\
\grammore{\term{endmodule}}

The optional ``\term{\# (} \many{\nterm{Ide}}{\term{,}} \term{)}'' on
the first line represents parameters to the module.  If omitted, it is
treated as \term{\#()}, the empty argument list.

In the abstract syntax, such a module definition is converted into a
\term{Binding} that binds the \nterm{Ide} (name of module definition)
to a lambda-expression whose arguments are the parameters and body is
a module-expression:

\begin{Verbatim}[frame=single, commandchars=\\\{\}]
{\it Binding:}    ({\it{}Ide}, (Lambda [{\it{}{Ide}} ... {\it{}{Ide}}] {\it{}{module-expression}}))
\end{Verbatim}

{\bf Abstract syntax} for module-expressions:

\begin{Verbatim}[frame=single, commandchars=\\\{\}]
data Expr = ...
          | ModuleExpr [Binding] [Rule] [MethDef]    -- {\it{}module-expression}
\end{Verbatim}

(\term{Expr} is defined in more detail later.)

% ----------------

\subsection{Rules}

Rules are defined inside modules.

{\bf Concrete syntax:}

\gram{Rule}{\term{rule} \nterm{Ide} \term{(} \opt{Expr} \term{)} \term{;}} \\
\grammore{\hm \nterm{Stmt}} \\
\grammore{\hm ...} \\
\grammore{\hm \nterm{Stmt}} \\
\grammore{\term{endrule}}

The \nterm{Expr} is the rule condition (a boolean expression).
``\many{\nterm{Stmt}}{\hm}'' is a list of zero or more statements
constituting the rule body.  In the abstract syntax, the rule
condition and body are combined into a ``When'' expression (described
in more detail later in Sec.~\ref{sec_syntax_expr}), and the
statements are combined into a Block.

\begin{Verbatim}[frame=single, commandchars=\\\{\}]
    Rule  {\it{}Ide}  (When {\it{}{Expr}} (Block [{\it{}{Stmt}} ... {\it{}{Stmt}}]))
\end{Verbatim}

{\bf Abstract syntax:}

\begin{Verbatim}[frame=single, commandchars=\\\{\}]
data Rule = Rule Ide Expr

data Expr = ...
          | Block [Stmt]            -- Block
          | When Expr Expr          -- When guard-cond value
\end{Verbatim}


% ----------------

\subsection{Method Definitions in Modules}

Methods are defined inside modules.

{\bf Concrete syntax:}

\gram{MethodDef}{\term{method} \nterm{MethodKind} \nterm{Ide}
                 \term{(} \many{\nterm{Ide}}{\term{,}} \term{)}         % args
                 \opt{\term{if} \term{(} \nterm{Expr} \term{)}}        % method condition
                 \term{;}} \\
\grammore{\hm \many{\nterm{Stmt}}{\hm}} \\
\grammore{\term{endmethod}}

\gram{MethodKind}{\term{V} \hm \alt \hm \term{A} \hm \alt \hm \term{AV}}

The \nterm{MethodKind} specifies whether this is a value method (V),
Action method (A) or ActionValue method (AV).  Value methods are pure
(no side-effects). Action methods just have a side effect, and do not
return any value. ActionValue methods have a side effect and return a
value.

This is followed by an \nterm{Ide}, the name of the method, and the
arguments of the methods in parentheses. This is optionally followed
by the keyword \term{if} and the method condition, a boolean
\nterm{Expr}. A missing method condition is treated as \verb|if (True)|.
Then we have the method body, which is a sequence of
statements.

In the abstract syntax, this is converted into a \term{MethDef} that
contains the method kind, the method name, and a When-expression.
whose condition is the method condition and whose body is a
lambda-expression.  The lambda-expression's arguments are the method
arguments, and the body is a Block containing the method's statements:

\begin{Verbatim}[frame=single, commandchars=\\\{\}]
    MethDef  \{methKind = "A/V/AV",
              methName = {\it{Ide}},
              methBody = When {\it{Expr}}
                             Lambda [{\it{Ide}} ... {\it{Ide}}]
                                 Block [{\it{Stmt}} ... {\it{Stmt}}]
             \}
\end{Verbatim}


{\bf Abstract syntax:}

\begin{Verbatim}[frame=single, commandchars=\\\{\}]
data MethDef = MethDef \{ methKind :: String    -- "V", "A" or "AV"
                       , methName :: Ide
                       , methBody :: Expr      -- When-expr
                       \}
\end{Verbatim}

% ----------------

\subsection{Expressions}

\label{sec_syntax_expr}

Finally, we can describe all the different kinds of expressions.

{\bf Concrete syntax:}

\gram{Expr}{\term{(} \nterm{Expr} \term{)}}                                               \hfill Parenthesized expression \\
\grammore{\nterm{Ide}}                                                                    \hfill Identifier \\
\grammore{\term{(} \term{)}}                                                              \hfill Void literal \\
\grammore{\nterm{IntegerLiteral}}                                                         \hfill integer literal \\
\grammore{\nterm{StringLiteral}}                                                          \hfill string literal \\
\grammore{\term{True} \alt \hm \term{False}}                                              \hfill boolean literal \\
\grammore{\term{if} \term{(} \nterm{Expr} \term{)} \nterm{Expr} \term{else} \nterm{Expr}} \hfill conditional  \\
\grammore{\term{while} \term{(} \nterm{Expr} \term{)} \nterm{Expr}}                       \hfill loop  \\
\grammore{\nterm{Expr} \term{.} \nterm{Ide}}                                              \hfill Method value \\
\grammore{\term{begin} \many{\nterm{Stmt}}{\hm} \term{end}}                               \hfill block

{\bf Abstract syntax:}

\begin{Verbatim}[frame=single, commandchars=\\\{\}]
data Expr = EIde String             -- Identifier

          | Void                    -- Constant void
          | ConstI Int              -- Constant integer
          | ConstS String           -- Constant string
          | PrimFn Opcode           -- Constant primitive function value
          | Lambda [Expr] Expr      -- args have form 'EIde x'
          | App [Expr]              -- apply e0 to e1, e2, ...
          | If Expr Expr Expr       -- Conditional
          | While Expr Expr         -- Loop

          | MethodVal Expr MethName -- Method value: Expr is a module instance

          | Block [Stmt]            -- Block

          | When Expr Expr          -- When guard-cond value

          | ModuleExpr [Binding] [Rule] [MethDef]

          -- The following only appears during static elaboration
          | ModuleInstRef Int   -- index into global IntMap of module instances

          -- The following only appear during dynamic execution
          | Unavail             -- due to a method condition being False
          | Closure Env Expr    -- value of a lambda

data Stmt = StmtExpr Expr
          | StmtBinding Binding
\end{Verbatim}

The first several alternatives are quite conventional (identifiers,
constants, lambda-expressions, applications, conditionals, loops).  As
usual, lambda-expressions will evaluate to closures at run time.  A
Block-expression contains a list (sequence) of statements, each of
which can be either a let-binding or an expression (typically a
method-call or function-call).  The final value of the
Block-expression is the value of the last item: Void if it is a
binding, and the expression value if it is an expression.

A When-expression has a guard (a boolean expression) and a body (an
expression).  At run time, a When-expression evaluates to the special
value \term{Unavail} if the guard expression evaluates to False or is
itself \term{Unavail}.  The body is only executed if the guard
evaluates to True.

During static elaboration, a \term{ModuleExpr} creates a new module
instance which has a unique integer identifier (more about this
later), and is replaced by a \term{ModuleInstRef~$j$} value that
refers to it.

A method value contains an expression and a method name.  At run time,
the expression evaluates to a particular module instance reference,
and the method name selects a method in that module instance.  In an
application, the function value may be a primitive function value, a
standard lambda-closure, or a method value.

% ================================================================

\section{Parsing from concrete syntax to abstract syntax}

\label{sec_parsing}

The Haskell file {\tt Parser.hs} defines a parser from concrete syntax
into abstract syntax, using the standard Haskell \term{Parsec}
library.  The main function is:

\begin{Verbatim}[frame=single, commandchars=\\\{\}]
type Program    = ([Binding], Schedule)

parseProgramFromString :: String -> Program
\end{Verbatim}

Its input is a String representing a complete program (typically taken
from a source file), and its output is a pair: the list of top-level
bindings (module definition bindings), and the schedule (list of
hierarchical names identifying rule instances).

This is a vanilla use of Haskell's Parsec library, and is not
particulary interesting for a discussion about semantics of the
abstract syntax, so we do not discuss it further here.

% ================================================================

\section{Static Elaboration}

\label{sec_static_elaboration}

The Haskell code for static elaboration is in the file \term{StaticElab.hs}.

The source program is parsed into a list of bindings and a schedule.
Each binding associates an identifier to a lambda-expression whose
body is a module-expression, as described in
Sec.~\ref{sec_syntax_module_definitions}.  The schedule is not
relevant for static elaboration; it will only play a role in dynamic
semantics, discussed in Sec.~\ref{sec_dynamic_semantics}.

One of the bindings in a source program must bind the distinguished
name \term{main} to module with no parameters (more specifically, to a
niladic lambda-expression whose body is a module-expression).  The
\emph{static elaboration} step applies this lambda expression (to zero
arguments) to produce a statically elaborated program, which is
(abstractly) a description of a concrete hardware module hierarchy
with behavior expressed as rules.

The process of static elaboration is recursive: applying \term{main}
will result in a module instance; the bindings in this module, when
evaluated, will create other module instances.  The bindings inside
those module instances may, in turn, bind further module instances,
and so on, recursively.  A particular source-program module definition
may be instantiated multiple times.  Ultimately we have a hierarchy
(i.e., a \emph{tree}) of module instances, with an instance of
\term{main} at the root.

% ----------------

\subsection{Module instances}

A module instance is either an instance of a primitive module
(built-in) or is an instance of a user-defined module in the source
program (the value of a module-expression):

\begin{Verbatim}[frame=single, commandchars=\\\{\}]
data ModuleInst = PMIReg      Int               -- reg data
                | PMICReg Int Int               -- # of ports, reg data
                ...
                {\it ... other primitives}
                ...
                | UMI     Env [Rule] [MethDef]  -- User-defined module instance
\end{Verbatim}
(Mnemonics: PMI = Primitive Module Instance, UMI = User Module Instance).

Here we show ``Registers'' and ``Concurrent Registers'' as primitives,
but these can be replaced or extended with an arbitrary set of
primitives such as FIFOs, RAMs, and so on.  Each primitive has its own
representation of internal state.  For example, the \term{Int} value
for Registers represents the current register data contents.  For
Concurrent Registers, the first \term{Int} represents the number of
its concurrent ports, and the second \term{Int} represents its current
data contents.

Static elaboration is a classical lambda-calculus ``Eval'' of the
source program.  This is supplied with an initial environment that
binds the names of constructors for primitive modules to primitive
functions:

\begin{Verbatim}[frame=single, commandchars=\\\{\}]
        primitivesEnv :: Env
        primitivesEnv = [ (EIde "mkReg",  PrimFn PrimMkReg)
                        , (EIde "mkCReg", PrimFn PrimMkCReg)
                        ] 
\end{Verbatim}

When applied to arguments (during static elaboration), these primitive
functions create (instantiate) \term{PMI} module instances.

When module-expressions are evaluated, they create \term{UMI} module
instances.  User-defined module instances are closures---the
\term{Env} component represents the bindings in the module, which may
be used in its rules and method definitions.

% ----------------

For the dynamic semantics the module hierarchy is not important: the
only important thing is that a rule in one module instance may invoke
a method in another module instance, and thus must have a way to refer
to another module instance.  Thus, during static elaboration, we give
each module instance a \emph{unique integer name}, and we collect all
module instances in a \emph{mapping} (a Haskell \term{IntMap} object)
that maps each unique integer name to a pair--- the module instance's
hierarchical name, and its instance value.  We call this mapping a
\term{System}:

\begin{Verbatim}[frame=single, commandchars=\\\{\}]
type System = Intmap (HierName, ModuleInst)
\end{Verbatim}

Each time we evaluate a ``\verb|ModuleExpr [Binding] [Rule] [MethDef]|''
expression, we create a new \term{ModuleInst} in the global
\term{system}, and replace the expression with a \verb|ModuleInstRef Int|
value, which is a reference to the module instance.

\begin{Verbatim}[frame=single, commandchars=\\\{\}]
data Expr = ...
          | ModuleInstRef Int    -- index into global IntMap of module instances
\end{Verbatim}

% ----------------

\subsection{The Static Elaboration function}

Static elaboration is carried out by this function:

\begin{Verbatim}[frame=single, commandchars=\\\{\}]
staticElab :: [Binding] -> (System, Expr)
staticElab bindings = eval  []  ["main"]  DI.empty  (App [EIde "main"])
    where
        ...
\end{Verbatim}

The argument is the list of bindings from the source file produced by
the parser.  It merely calls the \term{eval} function, to be described
next, giving it:
\begin{tightlist}
\item an initially empty environment: \term{[]},
\item an initial hierarchical name: \term{["main"]}
\item an initially empty map of module instances: \term{DI.empty},
\item and an initial application: \term{(App [EIde "main"])} i.e., \term{main} applied to zero arguments.
\end{tightlist}

Static elaboration is performed by a classical mutually recursive
Eval-Apply pair performing evaluation and application (locally inside
the \term{staticElab} function):

\begin{Verbatim}[frame=single, commandchars=\\\{\}]
  eval :: Env -> HierName -> System -> Expr -> (System, Expr)
  app :: HierName -> System -> Expr -> [Expr] -> (System, Expr)
\end{Verbatim}

In both \term{eval} and \term{app}:
\begin{tightlist}

\item The \term{HierName} argument is the current position in the module instance
hierarchy, i.e, the list of module instance names above this recursive point.

\item The \term{System} argument is the collection of module instances accumulated so far.

\end{tightlist}

and the result is a pair: the updated \term{System}, and an
\term{Expr} representing an evaluated value.  In \term{eval}:
\begin{tightlist}

\item The \term{Env} argument is the current environment binding identifiers to values.

\item The \term{Expr} argument is the expression being evaluated.

\end{tightlist}

Most of the clauses in \term{eval} are classical (identifiers,
constants, primitive functions, lambdas, applications).  The only
unusual clause is for module definitions:

\begin{Verbatim}[frame=single, commandchars=\\\{\}]
eval env hn system (ModuleExpr bs rs mds) =
                let
                    -- Convert the bindings into an environment,
                    -- sequentially evaluating the right-hand sides
                    f (system, env) (lhs @ (EIde s), rhs) = (system', env')
                        where
                            (system', v) = eval env (hn ++ [s]) system rhs
                            env' = env ++ [(lhs, v)]
                    (system', env') = foldl f (system, []) bs

                    -- Create a user-defined module instance
                    new_umi = UMI (env ++ env') rs mds
                in
                    extendSystem system' hn new_umi
\end{Verbatim}

A module-expression contains bindings \term{bs} of identifiers to
expressions (typically module instantiations), rules \term{rs} and
method definitions \term{mds}.  Here, we leave \term{rs} and
\term{mds} untouched.  We convert \term{bs} into an environment
\term{env'} by evaluating each right-hand side expression to a value.
We do this sequentially, using \term{foldl}, so that values from
earlier bindings are in scope for later bindings.  We create a new
user-defined module instance (\term{new\_umi}).  Finally,
\term{extendSystem} adds this module instance to the system,
allocating a new unique integer identifier $j$, returning the new
extended system and \term{ModuleInstRef~$j$} as the value of the
expression.

In \term{app}:

\begin{tightlist}

\item The \term{Expr} argument is the value of a function to be
applied (either a primitive, or the closure of a lambda expression if
user-defined).

\item The \term{[Expr]} argument is a list of argument values for the function.

\end{tightlist}

Most of the clauses in \term{app} are classical (primitive functions
and lambda-closures).  The only unusual clauses are for primitive
module constructors.  In each case, we simply extend the incoming
\term{system} with new primitive module instances:

\begin{Verbatim}[frame=single, commandchars=\\\{\}]
app hn system (PrimFn PrimMkReg) [ConstI j]
    = extendSystem system hn (PMIReg j)

app hn system (PrimFn PrimMkCReg) [ConstI n, ConstI j]
    = extendSystem system hn (PMICReg n j)
\end{Verbatim}

% ================================================================

\section{Dynamic Semantics}

\label{sec_dynamic_semantics}

Having parsed and elaborated a program, we are ready to turn to the
central focus of this document--the dynamic semantics of rules.

% ----------------------------------------------------------------

\subsection{The central intuitions behind clocks and schedules}

Hardware execution is a sequence of clocks.  In each clock we execute
a sequence of rules\footnote{We remind that the concept of a sequence
of rules executing in a clock is only a logical concept to explain
semantics and understand behavior.  It may not be discernible at all
in the clocked digital hardware that is the output of the {\it bsc}
compiler.}.  We call a sequence of rules a \emph{schedule}.

As explained in Sec.~\ref{sec_schedules}, the semantics takes a
schedule as a parameter, and takes no position on how that schedule is
constructed, it just explains correct execution of the given schedule.

Thus, the central question answered by the semantics is:
\begin{tabbing}
\hm \= In the current clock, \\
    \> given that we have already performed rule executions $r_0$, $r_1$, ..., $r_{j-1}$ so far in this clock, \\
    \> and the schedule has determined that we must next attempt rule $r_j$, \\
\\
    \> What is a correct execution of $r_j$?
\end{tabbing}
This step is repeated (within this clock) until the schedule offers no more rules for this clock.

And all this is repeated forever for multiple clocks.

Correct execution of a rule is straightforward, and is summarized
informally in Fig.~\ref{Fig_central_idea}.  Internalizing this will
make it easy to understand the formal description that follows.  This
is also (roughly) the mental model used every day by the practicing
BSV programmer to understand BSV program execution and the schedules
produced by the \emph{bsc} compiler.

\begin{figure}[htbp]
\fbox{
\begin{minipage}{\textwidth}

\begin{center}
  \begin{minipage}{\hlessmmmm}\bf
    \begin{center}
      Correct execution of a rule in a clock, given that we have previously
      executed certain rules in this clock.

      (Informal Description)
    \end{center}
  \end{minipage}
\end{center}

\begin{itemize}

\item We first ``evaluate'' the rule (which is a When-expression).
This is mostly a conventional lambda-calculus Eval-Apply pair, with a
few variations:

  \begin{itemize}

  \item In any When-expression (typically rule and method bodies), if
      the guard condition evaluates to \term{False}, the
      When-expression evaluates to the special value \term{Unavail}.
      In fact, the guard condition may itself evaluate to
      \term{Unavail}, in which case also the When-expression's value
      is \term{Unavail}.

  \item We collect a record of every method called, whether on a
     primitive module instance or on a user-defined module instance.

  \item We collect \emph{but do not perform} applications of primitive
     Action method values to their argument values.  Note, therefore,
     that, so far, \emph{rule evaluation has no side effects on the
     system state.}

  \end{itemize}

\item After rule evaluation, we decide whether the rule will fire or
not.  Any of the following conditions will stop a rule from firing.

  \begin{itemize}

  \item There is at least one \emph{intra-rule conflict}, i.e., the
      rule invokes two methods $m_1$ and $m_2$ that are not allowed to
      be invoked within the same Action
      (simultaneously/instantaneously/parallel).  These prohibitions
      are axiomatic properties of methods on primitive modules, and
      are given in a fixed intra-rule conflict table.

  \item There is at least one \emph{inter-rule conflict}, i.e., a rule
      executed earlier in this clock invokes a method $m_1$ and this
      rule invokes a method $m_2$, and they are not allowed to be
      executed in that order.  These prohibitions are axiomatic
      properties of methods on primitive modules, and are given in a
      fixed inter-rule conflict table.


  \item The rule evaluates to \term{Unavail}.  Either the rule
      condition expression, or the condition of some method (either in
      the rule condition expression or its body), evaluates to False.

  \end{itemize}

\item Finally, if we decide that the rule will fire, we execute the
  primitive Actions that we collected during rule evaluation.  This
  transforms the system state.  Conceptually, all these actions are
  performed simultaneously and instantaneously, i.e., in parallel.

\end{itemize}

\end{minipage}}

\caption{\label{Fig_central_idea}The central idea in rule semantics: correct rule execution}
\end{figure}

% ----------------

\subsection{Formal description}

The code for the dynamic semantics is in the file \term{Semantics.hs}.
We skip over the first few definitions, \verb|exec_n_clocks|,
\term{findRule} and \term{initEnv} because they are about specifying
schedules (here, we arbitrarily read a schedule at the end of each
source file) and repeated execution of a schedule across clocks, etc.,
which, as explained earlier, are not the central focus of this
document.

Instead, we pick it up at the following function:

\begin{Verbatim}[frame=single, commandchars=\\\{\}]
exec_1_clock :: System -> [(HierName,Env,Rule)] -> IO (Bool, System)
exec_1_clock system ruleSeq = do (anyfired, system', _)
                                     <- foldM f  (False, system, [])  ruleSeq
                                 return (anyfired, system')
    where
        -- mcsPrev is the collection of methods-called in rules
        -- previously fired in this clock
        f (anyfired, system, mcsPrev) rule =
            do (fired, system', mcsPrev') <- execRule system mcsPrev rule
               return (anyfired || fired, system', mcsPrev')
\end{Verbatim}

The \term{system} argument is the \term{IntMap} from integers to
module instances.

The \term{ruleSeq} argument is a list of rule instances representing
the schedule.  Each is a triple of a hierarchical name that uniquely
identifies the rule, an environment representing the bindings in the
module instance where this rule is found, and a
Rule-expression.\footnote{ There is no assumption that the schedule is
constant across clocks.  In fact, an alternative formulation may
replace the \term{ruleSeq} argument by an oracle function on the
current system state and list of rules fired so far in a clock, i.e.,
it can take a last-minute instantaneous decision on which rule should
be considered next.}.

The result of \term{exec\_1\_clock} is a pair. The \term{Bool} component
indicates whether any rule fired in this clock. The \term{System}
component is the updated system state (which will be the same as the
\term{system} argument if the rule did not fire).

The \verb|exec_1_clock| function is straightforward, just using
\term{foldM} to apply the \term{execRule} function over the rules in
the schedule, carrying along three values: a \term{Bool} indicating
whether any rule fired, the current system state, and \term{mcsPrev},
the collection of methods called in rules that have fired so far in
this clock.

A ``method-called'' is simply a method value, which is an expression:

\begin{Verbatim}[frame=single, commandchars=\\\{\}]
type MC = Expr    -- MethodVal (ModuleInstRef j) MethName
\end{Verbatim}

% ----------------

\subsubsection{Evaluating a Rule Expression}

The central function of this semantics is this one, which is a direct
implementation of the informal explanation of
Fig.~\ref{Fig_central_idea}.

\begin{Verbatim}[frame=single, commandchars=\\\{\}]
execRule :: System -> [MC] -> (HierName, Env, Rule) -> IO (Bool, System, [MC])
execRule  system  mcsPrev  (hname, env, Rule name e) =
  do -- Eval this rule
     (v,actions,mcsThis) <- evalExpr system env e

     let intra_cfs = intraRuleConflicts  system  mcsThis
     let inter_cfs = interRuleConflicts  system  mcsPrev  mcsThis
     let hw_cfs    = hwConflicts         system  mcsPrev  mcsThis

     -- has intra-rule method conflicts
     if not (intra_cfs == []) then return (False, system, mcsPrev)

     -- has inter-rule method conflicts
     else if not (inter_cfs == []) then return (False, system, mcsPrev)

     -- has hardware method conflicts
     else if not (hw_cfs == []) then return (False, system, mcsPrev)

     -- Rule is not enabled (rule cond, or some method cond, is false)
     else if (v == Unavail) then return (False, system, mcsPrev ++ mcsThis)

     -- Rule is enabled (rule-cond is True and all method-conds are True)
     -- Perform the rule's actions, updating the system state
     else do system' <- doActions system actions
             return (True, system', mcsPrev ++ mcsThis)
\end{Verbatim}

The \term{mcsPrev} argument is a collection of methods called in rules
that have previously fired in this clock.  The third argument is a
rule instance (its hierarchical name, an environment representing the
bindings in its surrounding module, and the Rule-expression itself).
The output is a boolean indicating whether the rule fired, the updated
system state (same as input system state if the rule did not fire),
and an updated collection of methods called (\term{mcsPrev} augmented
by methods called in this rule, if any).

The \term{evalExpr} function evaluates an expression to a value, which
is just an expression restricted to certain forms:

\begin{Verbatim}[frame=single, commandchars=\\\{\}]
type Val    = Expr    -- Unavail
                      -- Void, ConstI j, ConstS s,
                      -- PrimFn op,
                      -- Closure env (Lambda ...),
                      -- MethodVal (MmoduleInstRef j) MethName
\end{Verbatim}

The \term{evalExpr} function is a classical lambda-calculus Eval-Apply
pair, with a few variations.

\begin{Verbatim}[frame=single, commandchars=\\\{\}]
evalExpr :: System -> Env -> Expr -> IO (Val, [Action], [MC])
evalExpr system env e = eval e
    where
        eval :: Expr -> IO (Val, [Action], [MC])

        {\it ... clauses for classical eval of identifiers, constants, ...}
        {\it ...   primitive functions, lambdas ...}
        {\it ...   conditionals, loops, blocks ...}
        {\it ...}
\end{Verbatim}

We omit here the classical clauses (identifiers, constants,
conditionals, lambdas, ...).  If the expression is \verb|(MethodVal e methName)|,
then \term{e} evaluates to a module-instance reference of
the form \term{ModuleInstRef}~$j$ where $j$ is an index into the
\term{system} IntMap and selects a module instance.

\begin{Verbatim}[frame=single, commandchars=\\\{\}]
        eval (MethodVal e methName) =
            do (miref, [], mcs) <- evalExpr system env e
               return (MethodVal miref methName, [], mcs)
\end{Verbatim}

Since this evaluation cannot have side-effects, the pattern-match
asserts \term{[]} for the actions collected.

Evaluating a When-expression can return \term{True}, \term{False} or
\term{Unavail}.  The boolean values are encoded as \verb|ConstI 1| and
\verb|ConstI 0|.

\begin{Verbatim}[frame=single, commandchars=\\\{\}]
        eval (When econd ebody) =
            do (vcond, [], mcs) <- evalExpr system env econd
               if ((vcond == Unavail) || (vcond == (ConstI 0)))
               then return (Unavail, [], mcs)
               else do (vbody, as', mcs') <- evalExpr system env ebody
                       if (vbody == Unavail) then return (Unavail, [], mcs)
                       else return (vbody, as', mcs++mcs')
\end{Verbatim}

Evaluating an application is classical: evaluate the list of
expressions representing the function and the arguments, and then
apply the function value to the argument values except, if any of
those values are \term{Unavail}, just return \term{Unavail}.

\begin{Verbatim}[frame=single, commandchars=\\\{\}]
        eval (App es)   =
            do (vs, as, mcs) <- evalList es
               if any (== Unavail) vs then
                   return (Unavail, as, mcs)
               else do (v, as', mcs') <- apply system vs
                       return (v, as++as', mcs++mcs')
        ...
\end{Verbatim}

The \term{apply} function is classical for primitive functions and
closures.  For applying method values, we use the \term{applyMethod}
function.

\begin{Verbatim}[frame=single, commandchars=\\\{\}]
apply :: System -> [Expr] -> IO (Val, [Action], [MC])
apply system  (PrimFn opcode : vargs) =
    do (y, as) <- applyPrimFn (PrimFn opcode) vargs
       return (y, as, [])

apply system  (Closure env (Lambda xs e) : vargs) =
    evalExpr system (env ++ zip xs vargs) e

apply system  (mv @ (MethodVal miref methname): vargs) =
    applyMethod system mv vargs
\end{Verbatim}

The \term{applyPrimFn} function is just a direct implementation of
each primitive function's semantics.  We show the clause for \term{Plus} here:

\begin{Verbatim}[frame=single, commandchars=\\\{\}]
applyPrimFn :: Val -> [Val] -> IO (Val, [Action])
...
applyPrimFn (PrimFn Plus)   [ConstI x, ConstI y] = return (ConstI (x+y), [])
...
\end{Verbatim}

The \term{PrimDisplay} function is an Action, so we just collect it
for later execution after deciding that the rule will fire:

\begin{Verbatim}[frame=single, commandchars=\\\{\}]
...
applyPrimFn (PrimFn PrimDisplay) [ConstI x] =
    return (Void, [App [PrimFn PrimDisplay, ConstI x]])
applyPrimFn (PrimFn PrimDisplay) [ConstS s] =
    return (Void, [App [PrimFn PrimDisplay, ConstS s]])
...
\end{Verbatim}

The \term{applyMethod} function records the method called.  In
addition, for methods of primitive modules (like registers), it just
implements the primitive's semantics for value methods, or collects
the application for Action/ActionValue methods.

\begin{Verbatim}[frame=single, commandchars=\\\{\}]
applyMethod :: System -> Val -> [Val] -> IO (Val, [Action], [MC])
applyMethod system  mc@ (MethodVal  miref @ (ModuleInstRef j)  methname)  vargs =
    do let (hn, mi) = system IntMap.! j

       -- Record this method-called, whether it's a PMI or UMI
       let mcs = [mc]
       -- thisActions will only be used if this is not a Value method
       let thisActions = [App (mc:vargs)]

       case (mi, methname, vargs) of
           -- Ordinary registers
           (PMIReg v, "_read", []) -> return (ConstI v, [], mcs)
           (PMIReg v, "_write", [v']) -> return (Void, thisActions, mcs)
           ...
\end{Verbatim}

Application of user-defined methods is similar to application of
closures.  A user-defined method is in fact a closure---it is a
lambda-expression whose environment is the bindings in its enclosing
module instance. We evaluate the lambda-expression to a closure value,
and apply it to the arguments.

\begin{Verbatim}[frame=single, commandchars=\\\{\}]
           ...
           -- User-defined modules
           (UMI env rules methDefs, _, _) ->
               do -- Select the method by name, getting its body and kind
                  let [(mbody, mkind)] =
		    = [(mbody', mkind')| md <- methDefs
                                       , let MethDef mkind' methname' mbody' = md
                                       , methname == methname' ]
                  -- mbody is a lambda expr, evaluating to closure v
                  (v, actions, mcs') <- evalExpr system env mbody
                  -- let actions' = thisActions ++ actions
                  let actions' = if (mkind == "V") then []
                                 else (thisActions ++ actions)
                  if (v == Unavail) then return (Unavail, actions', mcs++mcs')
                  else do (v', actions'', mcs'') <- apply system (v:vargs)
                          return (v', actions'++actions'', mcs++mcs'++mcs'')
\end{Verbatim}

% ----------------

\subsubsection{Checking whether the rule has intra-rule conflicts}

After evaluating the rule (which has no side-effects, since we only
collect its Actions, we don't perform them yet), we check for
conflicts.  Checking whether the collection of methods-called within a
rule has any intra-rule conflicts is done with this function:

\begin{Verbatim}[frame=single, commandchars=\\\{\}]
intraRuleConflicts :: System -> [MC] -> [(MC,MC)]
intraRuleConflicts system mcs = ...
\end{Verbatim}

It checks each method-called in \term{mcs} with every other
method-called in \term{mcs}, and returns a list of all pairs that
conflict.  This check is just a test for presence in the fixed
\term{intraRuleConflictTable}.  An example entry in this table is
shown below.

\begin{Verbatim}[frame=single, commandchars=\\\{\}]
intraRuleConflictTable :: [ (ModuleType, MethName, MethName) ]
intraRuleConflictTable =
    [ -- Reg ----------------
      -- _write cannot be called more than once in a rule
      (Reg, "_write", "_write")
      ...
    ]
\end{Verbatim}

% ----------------

\subsubsection{Checking whether the rule has inter-rule conflicts}

Checking whether the collection of methods-called within a rule has
any inter-rule conflicts with any method called in rules previously
fired in this clock is done with this function:

\begin{Verbatim}[frame=single, commandchars=\\\{\}]
interRuleConflicts :: System -> [MC] -> [MC] -> [(MC,MC)]
interRuleConflicts  system  mcsPrev  mcsThis = ...
\end{Verbatim}

It checks each method-called in \term{mcsPrev} with each method-called
in \term{mcsThis}, and returns a list of all pairs that conflict.
This check is just a test for presence in the fixed
\term{interRuleConflictTable}.  An example entry in this table is
shown below.

\begin{Verbatim}[frame=single, commandchars=\\\{\}]
interRuleConflictTable :: [ (ModuleType, MethName, MethName) ]
interRuleConflictTable =
    [ -- Reg ----------------
      -- _write cannot precede _read
      (Reg, "_write", "_read")
      ...
    ]
\end{Verbatim}

% ----------------

\subsubsection{Checking whether the rule has hardware conflicts}

Checking whether there are any hardware conflicts is done with this function:

\begin{Verbatim}[frame=single, commandchars=\\\{\}]
hwConflicts :: System -> [MC] -> [MC] -> [(MC,MC)]
hwConflicts  system  mcsPrev  mcsThis = ...
\end{Verbatim}

It checks each method-called in \term{mcsThis} with each other
method-called in \term{mcsThis}, and it checks each method-called in
\term{mcsThis} with each method-called in \term{mcsPrev}.  Note that
hardware-conflicts between methods-called in \term{mcsPrev} would have
already been checked while executing earlier rules.  It returns a list
of all pairs that conflict.

A hardware conflict captures the fact that a hardware wire can only be
driven by one value in each clock.  Thus an Action or ActionValue
method always has a hardware conflict with itself (the ``enable''
input wire can only be driven once).  A Value method that has any
arguments always has a hardware conflict with itself (the argument
input wires can only be driven once).  To put it another way, only a
niladic Value method does not have a hardware conflict with itself.
This is captured in the following function.

\begin{Verbatim}[frame=single, commandchars=\\\{\}]
hwConflict :: System -> MC -> MC -> Bool
hwConflict system  (MethodVal (ModuleInstRef j)  methname)
                   (MethodVal (ModuleInstRef j2) methname2)
                       | j /= j2               = False    -- different modules
                       | methname /= methname2 = False    -- different methods
                       | otherwise             = b
  where
      (hname, mi) = (system IntMap.! j)
      b = case mi of
              (PMIReg _)    -> (methname == "_write") -- Action method

              {\it ... other primitives ...}

              (UMI _ _ _)   ->    -- Method on user-defined module
                  let
                      md               = select_UMI_method  methname  mi
                      Lambda args body = methBody md
                      niladic_Value_method = (   (methKind md == "V")
                                              && (length args == 0))
                  in
                      not niladic_Value_method
\end{Verbatim}

% ----------------

\subsubsection{Firing the rule (performing its actions)}

Finally, after evaluating a rule, checking that it has no inter-rule,
intra-rule or hardware conflicts, and checking that its rule and
method conditions are True (not \term{Unavail}), we can fire the rule,
i.e., perform its actions to update the system state.

The function \term{doActions} just uses \term{foldM} to iterate over
all the actions, repeatedly updating the system state:

\begin{Verbatim}[frame=single, commandchars=\\\{\}]
doActions :: System -> [Expr] -> IO System
doActions  system  actions = foldM  f  system  actions
    where f ...
\end{Verbatim}

Each action is executed by \term{doAction}.  Methods of primitive
modules just implement the primitive's semantics, updating the module
instance state to a new module instance state (like the register
\verb|_write| shown below).  Methods of user-defined modules are
ignored here: they were originally collected for checking hardware
conflicts, but they were also applied to their arguments to produce
(recursively), any underlying primitive-module actions.  Ultimately,
only primitive-module actions affect system state.

\begin{Verbatim}[frame=single, commandchars=\\\{\}]
doAction :: System -> Int -> MethName -> [Val] -> System
doAction  system  j  methName  args = system'
    where
        -- Get the primitive module instance from the system
        (hname, mi) = system IntMap.! j

        -- Upate the primitive module instance according to the method
        mi' = case (mi, methName, args) of
                  (PMIReg v,    "_write",  [ConstI v']) -> PMIReg v'
                  {\it ... other primitive-module actions ...}
                  {\it ... user-module actions are ignored ...}

        -- Put the updated primitive module instance back in the system
        system' = IntMap.insert j (hname, mi') system
\end{Verbatim}

And we're done with this rule (\term{execRule}). Function
\verb|exec_1_clock| will call \term{execRule} again for the next rule
in the schedule, until there are no more rules in the schedule.  Then,
\verb|exec_n_clocks| will call \verb|exec_1_clock| again for the next
clock, and so on.

% ================================================================

\section{Examples}

\label{sec_examples}

In this section we describe some examples that are in the
\term{Progs/} directory.  Each example is in a source file with
extension \term{.kbsv} (for kernel BSV).

% ----------------------------------------------------------------

\subsection{GCD}

It seems obligatory with hardware-design languages to start with a GCD
example (Greatest Common Divisor).  The code is in the file
\term{GCD.kbsv}. Here is the source code for the GCD module.

\begin{Verbatim}[frame=single, numbers=left, commandchars=\\\{\}]
module mkGCD;
    let  zero = 0;
    let x = mkReg (zero);
    let y = mkReg (zero);
    let busy = mkReg (zero);

  rules
    rule swap (   x._read()     >  y._read()
               && busy._read () == 1
               && y._read ()    != 0);
        x._write (y._read ());
        y._write (x._read ())
    endrule

    rule subtract (   x._read ()    <= y._read ()
                   && busy._read () == 1
                   && y._read ()    != 0);
        y._write (y._read () - x._read ())
    endrule

  methods
    method A start (num1, num2) if (busy._read () == 0);
        x._write (num1);
        y._write (num2);
        busy._write (1)
    endmethod

    method AV getResult () if ((busy._read () == 1) && (y._read () == 0));
        busy._write (0);
        x._read ()
    endmethod
endmodule
\end{Verbatim}

L.2 (line 2) is just a trivial binding giving the name \term{zero} to
the constant integer 0, just to demonstrate that you can have
conventional let-bindings of names to expressions.  The next three
lines instantiate three registers with initial (reset) value 0, and
bind \term{x}, \term{y} and \term{busy} to those instances.

L.8-L.13 define the \term{swap} rule.  The rule condition tests if
$x>y$, if we are busy and if $y\neq 0$.  Note, we're using the
\verb|_read()| method on a register to retrieve its value (in BSV
these method calls can be omitted and \emph{bsc} automatically fills
them in).  The body of the rule, L.11-L.12 effectively swap the $x$
and $y$ values.  Here we use the \verb|_write()| method to store a
value into a register (in BSV we typically use the more convenient
infix \verb|<=| syntax).

L.15-L.19 define the \term{subtract} rule. Here the rule condition
tests if $x\leq y$ (and busy and $y\neq 0$).  The rule body assigns
$y-x$ to $y$.

L.22-L26 define the \term{start} method.  The method condition ensures
it can only be invoked when we are not busy.  The method body stores
the arguments into $x$ and $y$, and store True (coded as integer 1)
into \term{busy}.

L.28-L.32 define the \term{getResult} method. The method condition
ensures it can only be invoked when we are busy and $y=0$.  It returns
the value of $x$ and writes False (0) into \term{busy}, so that the
\term{start} method can now be invoked again.

Here is a ``\term{main}'' module to test our GCD module.

\begin{Verbatim}[frame=single, numbers=left, commandchars=\\\{\}]
module main;
    let state = mkReg (0);
    let gcd = mkGCD ();

  rules
    rule init (state._read () == 0);
        gcd.start (24, 16);
        state._write (1)
    endrule

    rule finish (state._read () == 1);
        let z = gcd.getResult ();
        $display ("The GCD is ");
        $display (z);
        state._write (2)
    endrule
  methods
endmodule
\end{Verbatim}

It instantiates a \term{state} register initialized to 0, and it
instantiates the GCD module.  Rule \term{init} can fire when
\term{state} is 0; it invokes \term{gcd.start(24,16)} and writes 1 to
\term{state}.  Rule \term{finish} can fire when \term{state} is 1; it
retrieves a result from \term{gcd}, displays it, and writes 2 to
\term{state}.  When \term{state} is 2 no rules will enabled, and this
will terminate the simulation.

The final piece of source text is a schedule:

\begin{Verbatim}[frame=single, numbers=left, commandchars=\\\{\}]
schedule
  [main, init]
  [main, finish]
  [main, gcd, swap]
  [main, gcd, subtract]
\end{Verbatim}

Each line is a hierarchical name.  For example, the last line
represents the \term{subtract} rule inside the \term{gcd} instance
inside the \term{main} instance.  In this example, the choice of
schedule is not important, since the rule conditions ensure that only
one of them can fire in a clock anyway.

% ----------------

\subsubsection{Executing the semantics}

We execute the semantics as follows:
\begin{Verbatim}[frame=single, numbers=left, commandchars=\\\{\}]
$ runghc  Main  100  Progs/GCD.kbsv
\end{Verbatim}

(This command is also available in the \term{Makefile}.)
``\term{runghc}'' is the standard Glasgow Haskell Compiler command to
execute a Haskell program.  ``\term{Main}'' refers to the file
\term{Main.hs} which is the top-level of the Haskell code provided,
which invokes all the parsing, static elaboration and semantic
functions discussed in this document.  ``\term{100}'' limits the
execution to a maximum of 100 (simulated) clock cycles.  Finally, we
provide the source file name.

The variable \term{debugLevel} in the file \term{Utils.hs} controls
the verbosity of the output.  When set to 0, it should only produce
the output of \verb|$display| functions.  We set it to 1 and captured
the output in the file \verb|log_GCD|.  It consists of three sections:
a rendering of the abstract syntax tree (AST) produced by parsing, a
rendering of the elaborated program produced by static elaboration
and, finally, a clock-by-clock trace of the execution.  We skip the
AST here, since it is not very interesting.

The first part of the elaborated program looks like this:

\begin{Verbatim}[frame=single, numbers=left, commandchars=\\\{\}]
---- Elaborated program
  0 ["main","state"]    PMIReg 0
  1 ["main","gcd","x"]    PMIReg 0
  2 ["main","gcd","y"]    PMIReg 0
  3 ["main","gcd","busy"]    PMIReg 0
  4 ["main","gcd"]
    UMI 
        zero = 0
        x = ModuleInstRef 1
        y = ModuleInstRef 2
        busy = ModuleInstRef 3
      RULES
        rule swap when (And(And(Gt(x._read(),y._read()),Eq(busy._read(),1)),
                        Neq(y._read(),0)));
          begin
            x._write(y._read());
            y._write(x._read());
          end
        rule subtract when (And(And(Leq(x._read(),y._read()),Eq(busy._read(),1)),
                                Neq(y._read(),0)));
          begin
            y._write(Minus(y._read(),x._read()));
          end
      METHODS
        method A start when (Eq(busy._read(),0));
          (lambda (num1,num2) 
            begin
              x._write(num1);
              y._write(num2);
              busy._write(1);
            end)
        endmethod
        method AV getResult when (And(Eq(busy._read(),1),Eq(y._read(),0)));
          (lambda () 
            begin
              busy._write(0);
              x._read();
            end)
        endmethod
    EUMI
\end{Verbatim}

This is the Haskell \term{IntMap} that maps integers (unique names) to
module instances.  For example, ``3'' refers to the register primitive
module instance (PMI) with hierarchical name \term{main.gcd.busy},
i.e., the \term{busy} register inside the \term{gcd} module instance
which is inside the \term{main} module instance.  The initial
(reset-value) data in that register is 0.

``4'' refers to the GCD user-defined module instance (UMI). In the
module, you can see that \term{busy} is bound to the value
\verb|ModuleInstRef 3|, i.e., the register instance discussed in the
previous para.  If we had two instances of the GCD module, each would
have its own entry in the IntMap.  Inside, the value bound to the name
\term{busy} would different in the two GCD instances, each referring
to its own register PMI, its own copy of \term{busy}.  The rules and
method definitions in the UMI are the same as in the original AST
(this is true in all UMIs---only the bindings are evaluated, the rules
and method definitions are left untouched).

The rest of the elaborated program is the UMI for \term{main}, and the
finally we have the value of the expression ``\term{main()}'' which is
the root of static elaboration.

\begin{Verbatim}[frame=single, numbers=left, commandchars=\\\{\}]
  5 ["main"]
    UMI 
        state = ModuleInstRef 0
        gcd = ModuleInstRef 4
      RULES
        rule init when (Eq(state._read(),0));
          begin
            gcd.start(24,16);
            state._write(1);
          end
        rule finish when (Eq(state._read(),1));
          begin
            let z = gcd.getResult()
            $display("The GCD is ");
            $display(z);
            state._write(2);
          end
      METHODS
    EUMI
value of main(): ModuleInstRef 5
\end{Verbatim}

After static elaboration, we see a clock-by-clock trace of execution.
In each clock, it says what happened with each rule in the schedule.
For example, clock 0:

\begin{Verbatim}[frame=single, numbers=left, commandchars=\\\{\}]
.... Clock 0
  Rule ["main","init"] fired
  Rule ["main","finish"]
    inter-rule conflicts:
      ["main","state"] _write > _read
  Rule ["main","gcd","swap"]
    inter-rule conflicts:
      ["main","gcd","x"] _write > _read
      ["main","gcd","x"] _write > _read
      ["main","gcd","y"] _write > _read
      ["main","gcd","y"] _write > _read
      ["main","gcd","y"] _write > _read
      ["main","gcd","busy"] _write > _read
  Rule ["main","gcd","subtract"]
    inter-rule conflicts:
      ["main","gcd","x"] _write > _read
      ["main","gcd","y"] _write > _read
      ["main","gcd","y"] _write > _read
      ["main","gcd","busy"] _write > _read
\end{Verbatim}

This shows that the first rule in the schedule, \term{main.init},
fired.  The second rule in the schedule, \term{main.finish} did not
fire: it had an inter-rule conflict, where it was trying to execute
the method \verb|main.state._read|, whereas rule \verb|main.init| has
already fired and executed \verb|main.state._write|, and \verb|_read|
cannot follow \verb|_write| in a clock.  And so on, for the remaining
rules in the schedule, \verb|main.gcd.swap| and
\verb|main.gcd.subtract|.

When \term{main.init} fires, it invokes the
\term{main.gcd.start(24,16)} method which stores 24 and 16 into
registers \term{main.gcd.x} and \term{main.gcd.y}, respectively.
Continuing through clock 5, we see:

\begin{Verbatim}[frame=single, numbers=left, commandchars=\\\{\}]
.... Clock 1
  ...
  Rule ["main","gcd","swap"] fired               {\it So: x = 16, y = 24}
.... Clock 2
  ...
  Rule ["main","gcd","subtract"] fired           {\it So: x = 16, y = 8}
.... Clock 3
  ...
  Rule ["main","gcd","swap"] fired               {\it So: x = 8, y = 16}
  ...
.... Clock 4
  ...
  Rule ["main","gcd","subtract"] fired           {\it So: x = 8, y = 8}
.... Clock 5
  ...
  Rule ["main","gcd","subtract"] fired           {\it So: x = 8, y = 0}
\end{Verbatim}

At this point, rule \term{main.finish} is enabled: its rule condition
is true since \term{main.state}=1, and the condition of the method it
calls, \term{main.gcd.getResult}, is true because
\term{main.gcd.busy}=1 and \term{main.gcd.y}=0.  Indeed, if we
continue to clock 6 we see:

\begin{Verbatim}[frame=single, numbers=left, commandchars=\\\{\}]
.... Clock 6
  ...
  Rule ["main","finish"] fired
$DISPLAY: The GCD is 
$DISPLAY: 8
  ...
\end{Verbatim}

and we see the outputs from the rule's \verb|$display| statements.

In clock 7, we see that no rules fire (they are all \term{Unavail}), so
the simulation prints the final system state and stops.  You can see
that the \term{main.state} register has value 2.
the \term{main.gcd.x} register has value 8,
the \term{main.gcd.y} register has value 0,
and the \term{main.gcd.busy} register has value 0.

\begin{Verbatim}[frame=single, numbers=left, commandchars=\\\{\}]
.... Clock 7
  Rule ["main","init"] unavail
  Rule ["main","finish"] unavail
  Rule ["main","gcd","swap"] unavail
  Rule ["main","gcd","subtract"] unavail
No rule fired on cycle 7
---- final system: 
  0 ["main","state"]    PMIReg 2
  1 ["main","gcd","x"]    PMIReg 8
  2 ["main","gcd","y"]    PMIReg 0
  3 ["main","gcd","busy"]    PMIReg 0
  4 ["main","gcd"]UMI ...
  5 ["main"]UMI ...
\end{Verbatim}

% ----------------------------------------------------------------

\subsection{PipelineFIFO}

The GCD example is nice, but it does not have any interesting
concurrency since only one rule condition is enabled at any time (this
is why the particular choice of its schedule does not matter either).
We now look at a more interesting example, a 2-stage pipeline between
which is a ``PipelineFIFO''.  The code is in the file
\term{Progs/PipelineFIFO.kbsv}.

The first module in the file is \term{mkPipelineFIFO}:

\begin{Verbatim}[frame=single, numbers=left, commandchars=\\\{\}]
module mkPipelineFIFO;
    let full = mkCReg (2, 0);
    let data = mkCReg (2, 0);

  rules
    -- no rules

  methods
    method A enq (x) if (full._read1 () == 0);
      data._write1 (x);
      full._write1 (1)
    endmethod

    method V notEmpty ();
      full._read0 ()
    endmethod

    method V first () if (full._read0 () == 1);
      data._read0 ()
    endmethod

    method A deq () if (full._read0 () == 1);
      full._write0 (0)
    endmethod
endmodule
\end{Verbatim}

This implements a one-element FIFO.  The \term{full} register is 1 if
the FIFO contains a datum, and 0 if it does not.  The \term{data}
register holds the datum itself.  The module contains no rules, just
method definitions.  The \term{enq} method may be called when
\term{full}=0, i.e., the FIFO is empty, in which case it stores the
datum \term{x} in the \term{data} register and stores 1 in
\term{full}.  The \term{notEmpty} method just returns the \term{full}
value. The \term{first} and \term{deq} methods may be invoked when
\term{full}=1.  The \term{first} method returns the datum in the FIFO,
and the \term{deq} method writes 0 \term{full}, i.e., it marks the
FIFO as empty.

If \term{full} were an ordinary register, we could never invoke
\term{enq} and \term{deq} in the same cycle. Both methods execute
\verb|full._read()| and write \verb|full._write()|.  Whether we
scheduled (a rule that invoked) \term{enq} before (a rule that
invoked) \term{deq} or vice versa, we would run into an inter-rule
conflict because \verb|_read()| must precede \verb|_write| in a clock.

Thus, we use \emph{Concurrent Registers} (CRegs) for \term{full} and
\term{data}.  The methods on the output side of the FIFO,
\term{notEmpty}, \term{first} and \term{deq}, use \verb|_read0| and
\verb|_write0|.  The method on the input side of the FIFO, \term{enq},
uses \verb|_read1| and \verb|_write1|.  The
\term{interRuleConflictTable} permits \verb|_write0| and \verb|_read0|
to precede \verb|_write1| and \verb|_read1| in a clock.  Thus,
provided the output side methods are scheduled before the input side
method, they can execute in the same clock.  Conceptually, we can
dequeue an element from the FIFO, thus emptying it, and enqueue a new
element into the FIFO, thus re-filling it, \emph{all within the same
clock}.

The main test program is shown here:

\begin{Verbatim}[frame=single, numbers=left, commandchars=\\\{\}]
module main;
    let x = mkReg (0);
    let f = mkPipelineFIFO ();

  rules
    rule feed;
      f.enq (x._read());
      x._write (x._read() + 1)
    endrule

    rule drain (f.notEmpty ());
      $display ("RESULT");
      $display (f.first ());
      f.deq ()
    endrule
  methods
endmodule
\end{Verbatim}

The \term{feed} rule repeatedly enqueues \term{x} into the FIFO and increments \term{x}.
The \term{drain} rule repeatedly dequeues an item from the FIFO and displays it.
At the bottom of the file, we show two candidate schedules (one of them is commented out):

\begin{Verbatim}[frame=single, numbers=left, commandchars=\\\{\}]
-- Sched 1 (pipelining happens)
schedule
  [ main, drain ]
  [ main, feed ]

/*
-- Sched 2 (pipelining does not happen, rules alternate)
schedule
  [ main, feed ]
  [ main, drain ]
*/
\end{Verbatim}

% ----------------

\subsubsection{Executing the semantics with schedule 1}

We can execute the program; the transcript using Sched 1 is captured
in \verb|log_PipelineFIFO|.  Looking at the first 3 clocks, we see:

\begin{Verbatim}[frame=single, numbers=left, commandchars=\\\{\}]
.... Clock 0
  Rule ["main","drain"] unavail
  Rule ["main","feed"] fired
.... Clock 1
  Rule ["main","drain"] fired
$DISPLAY: RESULT
$DISPLAY: 0
  Rule ["main","feed"] fired
.... Clock 2
  Rule ["main","drain"] fired
$DISPLAY: RESULT
$DISPLAY: 1
  Rule ["main","feed"] fired
.... Clock 3
  Rule ["main","drain"] fired
$DISPLAY: RESULT
$DISPLAY: 2
  Rule ["main","feed"] fired
\end{Verbatim}

In clock 0, rule \term{drain} did not fire, since the FIFO is
initially empty, and so the \term{notEmpty}, \term{first} and
\term{deq} are not enabled. Rule \term{feed} fired, since the
\term{enq} method is enabled when the FIFO is empty.

In each subsequent clock, now that the FIFO is no longer empty, we see
that both rules fire, as expected.  First \term{drain} fires and
displays the result; this leaves the FIFO empty, enabling the
\term{feed} rule, which also fires.  The displayed outputs are 0, 1,
2, ... as expected.

The system stops after 100 clocks, as requested on the command line.
The final output is 99, as expected, and the final state can also be
seen:

\begin{Verbatim}[frame=single, numbers=left, commandchars=\\\{\}]
.... Clock 100
  Rule ["main","drain"] fired
$DISPLAY: RESULT
$DISPLAY: 99
  Rule ["main","feed"] fired
Cycle limit reached: 100
---- final system: 
  0 ["main","x"]    PMIReg 101
  1 ["main","f","full"]    PMICReg 2 1
  2 ["main","f","data"]    PMICReg 2 100
  3 ["main","f"]UMI ...
  4 ["main"]UMI ...
\end{Verbatim}

% ----------------

\subsubsection{Executing the semantics with schedule 2}

If we instead execute the program with the other schedule, Sched 2, we
get a different execution.  Its transcript is in file
\verb|log_PipelineFIFO_Sched2|.  The execution trace begins at clock
0:

\begin{Verbatim}[frame=single, numbers=left, commandchars=\\\{\}]
.... Clock 0
  Rule ["main","feed"] fired
  Rule ["main","drain"]            {\it did not fire}
    inter-rule conflicts:
      ["main","f","data"] _write1 > _read0
      ["main","f","full"] _write1 > _read0
      ["main","f","full"] _write1 > _read0
      ["main","f","full"] _write1 > _read0
      ["main","f","full"] _write1 > _write0
\end{Verbatim}

In clock 0, the FIFO is empty, so rule \term{feed} is enabled and
fires, ultimately using methods \verb|full._read1()|,
\verb|full._write1()| and \verb|data._write1()|.  Then, rule
\term{drain} cannot fire, since it ultimately invokes
\verb|full._read0()|, \verb|full._write0()| and \verb|data._read0()|,
and these cannot follow the already-invoked \verb|full._read1()|,
\verb|full._write1()| and \verb|data._write1()| (inter-rule conflict).
Moving on to clock 1:

\begin{Verbatim}[frame=single, numbers=left, commandchars=\\\{\}]
.... Clock 1
  Rule ["main","feed"] unavail     {\it did not fire}
  Rule ["main","drain"] fired
$DISPLAY: RESULT
$DISPLAY: 0
\end{Verbatim}

In clock 1, rule \term{feed} cannot fire because the FIFO is full and
so the \term{enq} method is disabled. Rule \term{drain} can fire; it empties the FIFO, and
displays the dequeued item.

\begin{Verbatim}[frame=single, numbers=left, commandchars=\\\{\}]
.... Clock 2
  Rule ["main","feed"] fired
  Rule ["main","drain"]            {\it did not fire}
    ...
\end{Verbatim}

In clock 2, rule \term{feed} is once again enabled since the FIFO is
empty, but this again prevents rule \term{drain} from firing, just
like in clock 0.  In subsequent clocks we see this pattern repeating,
with either rule \term{feed} or rule \term{drain} firing.

\begin{Verbatim}[frame=single, numbers=left, commandchars=\\\{\}]
.... Clock 3
  Rule ["main","feed"] unavail     {\it did not fire}
  Rule ["main","drain"] fired
$DISPLAY: RESULT
$DISPLAY: 1
.... Clock 4
  Rule ["main","feed"] fired
  Rule ["main","drain"]            {\it did not fire}
    ...
.... Clock 5
  Rule ["main","feed"] unavail
  Rule ["main","drain"] fired
$DISPLAY: RESULT
$DISPLAY: 2
.... Clock 6
  Rule ["main","feed"] fired
  Rule ["main","drain"]            {\it did not fire}
    ...
\end{Verbatim}

After we reach the requested cycle limit (100), simulation stops.  The
final output is 49 and the final state of register \term{x} is 51
(these were 99 and 101, respectively, with Sched 1).

\begin{Verbatim}[frame=single, numbers=left, commandchars=\\\{\}]
.... Clock 99
  Rule ["main","feed"] unavail
  Rule ["main","drain"] fired
$DISPLAY: RESULT
$DISPLAY: 49
.... Clock 100
  Rule ["main","feed"] fired
  Rule ["main","drain"]
    inter-rule conflicts:
      ...
Cycle limit reached: 100
---- final system: 
  0 ["main","x"]    PMIReg 51
  1 ["main","f","full"]    PMICReg 2 1
  2 ["main","f","data"]    PMICReg 2 50
  3 ["main","f"]UMI ...
  4 ["main"]UMI ...
\end{Verbatim}

In summary, observe that \emph{both schedules provide functionally
correct execution}: we enqueue and dequeue correct values.  In this
functional sense, neither schedule is ``correct'' or ``incorrect''.

However, Sched 1 admits more concurrency than Sched 2, because its
sequence is ``better aligned'' with the inter-rule ordering
constraints on the methods of the \term{data} and \term{full} CRegs.
In fact, the \emph{bsc} compiler will automatically pick Sched 1 using
this criterion, i.e., it tries to maximize concurrency.

% ----------------------------------------------------------------

\subsection{BypassFIFO}

The next example is the ``dual'' of the PipelineFIFO and is called the
BypassFIFO.  The code is in the file \term{Progs/BypassFIFO.kbsv}.  We
leave it as an exercise for the reader, with just the following
observations.

In the \term{PipelineFIFO}, the output-side methods (\term{notEmpty},
\term{first} and \term{deq}) of the FIFO used port 0 methods
(\verb|_read0| and \verb|_write0|) of the \term{full} and \term{data}
Concurrent Registers, and the input-side method (\term{enq}) used port
1 methods (\verb|_read1| and \verb|_write1|).  This allowed concurrent
input and output on the FIFO, provided the output side was scheduled
before the input side.

In the \term{BypassFIFO}, we do exactly the opposite: the input-side
method (\term{enq}) uses port 0 methods (\verb|_read0| and
\verb|_write0|), and the output-side methods (\term{notEmpty},
\term{first} and \term{deq}) use port 1 methods (\verb|_read1| and
\verb|_write1|) of the \term{full} and \term{data} Concurrent Registers.  This
allows concurrent input and output on the FIFO, provided the input
side is scheduled before the output side.

The main test program and the two candidate schedules are identical to
those in \term{PipelineFIFO.kbsv}.  The transcript using Sched 1 is
shown in file \verb|log_BypassFIFO_Sched1|, and using Sched 2 is shown
in file \verb|log_BypassFIFO|.  As with the PipelineFIFO, both show
functionally correct execution, just different levels of concurrency.
With Sched 2, both rules can fire on every clock; with Sched 1 they
alternate.

% ----------------------------------------------------------------

\subsection{PipeFwdBwd}

Our final example combines the last two and is representative of
structures we find, for example, in CPU pipelines.  In a CPU pipeline,
the main data flow is in the ``forward'' direction through the pipe,
and the pipeline stages are separated by PiplineFIFOs.  Thus, in each
clock, rules corresponding to all stages can fire concurrently,
allowing the entire pipeline to advance in each clock.  However,
within each clock, semantically the last stage executes first, then
the previous one, then its previous one, and so on until the first
stage.

Most CPUs also also have a ``reverse'' data flow, where later stages
may provide updated register values or branch-redirection information
back to early stages.  For consistent rule ordering with the forward
direction, therefore, we use BypassFIFOs in the reverse direction.

Our example is in the file \term{PipeFwdBwd.kbsv}.  It also has two
stages, using a PipelineFIFO in the forward direction from rule
\term{feed} to rule \term{drain}, and a BypassFIFO in the reverse
direction from rule \term{drain} to rule \term{feed}.

We leave it as an exercise to the reader to run \term{PipeFwdBwd.kbsv}
with each of the two schedules, and to explain their respective
behaviors.  For convenience, \verb|log_PipeFwdBwd| is a transcript of
execution with Sched 1, and \verb|log_PipeFwdBwd_Sched2| is a
transcript of execution with Sched 2.

% ----------------------------------------------------------------

% ================================================================

\section{Conclusion and closing observations}

\label{sec_conclusion}

We have described the dynamic execution semantics of BSV programs
formally by providing a Haskell program directly implementing the
semantics.

We also described a process of static elaboration, but that is not
central here; any alternative mechanism is fine that results in a
collection of module instances, where each module is either primitive
or is user-defined, where user-defined modules have rules and methods
that can invoke methods of primitives or other modules.

% ----------------------------------------------------------------

\subsection{Relationship to RTL and Chisel: Atomic transactions}

In the literature and training materials on BSV, we often refer to
rules as \emph{atomic transactions}.  Atomicity is trivial in the
semantics described in this document because we define \emph{serial}
execution of rules according to a schedule, and so there is no
question of interleaving of rule actions.

A central issue encountered in concurrent programming is that atomic
operations are non-modular, or non-compositional \cite{Harris2005}.
In other words, a system of modules, each of which implements
(locally) atomic operations does not guarantee that system-level
actions are atomic.  The classic example is a pair of shared bank
account modules with \term{deposit} and \term{withdraw} methods, each
of which are guaranteed atomic.  Two concurrent processes that attempt
a logical \term{transfer} operation by withdrawing from one account
and depositing into the other may nevertheless see inconsistent
results due to bad interleavings of these methods.

This is the reason why user-level atomic transactions cannot be
implemented correctly merely with libraries---the compiler and/or
run-time system must have visibility over \emph{all} actions of a
user-defined atomic transaction, and this set may span many
modules/functions/procedures (cf. ``Transactional Memory''
\cite{Larus2008a}).

This explains the fundamental power of BSV over traditional RTL
languages like Verilog~\cite{IEEEVerilog2001a},
SystemVerilog~\cite{IEEESystemVerilog2012a} and
VHDL~\cite{IEEEVHDL2002}.  The semantics of those languages is based
entirely on clocks and classical clocked digital logic.  There is no
concept of user-defined atomic transactions like rules in BSV.

The more recent language Chisel~\cite{BachRach2012} is a much
higher-level language than those Verilog, SystemVerilog or VHDL, but
its semantics is also fundamentally clock based, just like RTL.
Chisel has a syntactic ``when'' construct that at first glance look
like rule and method conditions in BSV, but they have nothing to do
with atomic transactions; they are just a syntactic means of
prioritizing updates to a state element, \emph{locally to a syntatic
scope} (in particular, they are not modular or compositional).

% ----------------------------------------------------------------

\subsection{Relationship to HLS (High Level Synthesis)}

The last 10 years has seen the emergence of a number of commercial
tools for hardware design under the category of ``High Level
Synthesis'', or HLS for short \cite{Coussy2009} (we'll refer to this
as ``classical HLS'').  The input design language for these tools is
an existing language, typically C or C++ (or SystemC, which is just
a C++ library \cite{IEEESystemC2011a}).  This is paradoxical, since C
and C++ are, formally, completely sequential languages, from which
classical HLS tries to generate hardware which, quintessentially, has
massive, fine-grained parallelism.  Classical HLS tools achieve this
with ``automatic parallelization'', which builds on several decades of
research into automatic parallelization of software for vector
computers and SIMD computers.  The central activity is to analyze
for-loops over dense rectangular arrays and to recognize when
iterations can be performed safely in parallel, from which they can be
mapped into parallel hardware.  As a consequence, classical HLS has
the same strengths and limitations as classical automatic
parallelization; a deeper discussion of these limitations may be found
in \cite{Edwards2005a}.

It should be evident that BSV and BSV compilation has almost nothing
in common with classical HLS.  The BSV programmer (unlike the C/C++
programmer) thinks in parallel, and explicitly designs and implements
parallel algorithms.  The ``high-level synthesis'' in BSV and the
\emph{bsc} compiler refers to a quite different set of properties and
activities:

\begin{itemize}

\item its extremely powerful Haskell-like type system,

\item its extremely powerful Haskell-like static-elaboration,

\item its automatic analysis of rules to produce highly concurrent rule schedules,

\item and its generation of efficient hardware from this analysis to
execute rules concurrently.

\end{itemize}

Unlike classical HLS, where hardware quality drops off sharply (and
may even becomes infeasible) as one moves out of the ``sweet spot'' of
simple loops on dense rectangular arrays, BSV has universal
applicability to all kinds of clocked digital hardware and is, in that
sense, a genuinely full replacement for RTL.

% ----------------------------------------------------------------

\subsection{Relationship of this formal semantics to the full BSV language}

Ideally, formal semantics should be defined directly on the source
language used by the programmer.  If this can be done, and the
semantic definition is ``simple'', then the everyday practicing
programmer can directly use it as a mental model while creating,
debugging and analyzing actual programs.  Everyday practicing
programmers find it difficult to usefully exploit formal semantics
that is defined on a kernel language that is too far removed from the
source language, or that involve esoteric source-to-source
transformations.

We have attempted to approach this ideal in this semantics.  The
language defined in Sec.~\ref{sec_syntax} is admittedly much smaller
and much simplified compared to full BSV.  But the principal
difference is in types and type-checking, and in a few syntatic
facilities like interface definitions and static-elaboration loops.
These are quite orthogonal to dynamic semantics and do not affect the
mental model held by the everyday practicing programmer.  We have also
omitted anything related to BSV packages, which again are just about
namespace management and separate compilation, and quite orthogonal to
dynamic semantics.

Thus, our kernel language captures the essence of BSV for purposes of
explaining dynamic semantics.  This semantics is in fact the mental
model one uses while doing BSV programming.

% ----------------------------------------------------------------

\subsection{Relationship of this formal semantics to the {\it bsc} compiler}

This formal semantics is purely a dynamic semantics explained directly
on the abstract syntax trees of source code---it takes no position on
how to find schedules (rule sequences), nor on compilation to actual
hardware.  As a result, this semantics is more liberal than what can
actually be achieved by any compiler (which, by definition, is working
statically).

The \emph{bsc} compiler only produces a \emph{static} schedule---a
single fixed schedule that is constant across clocks.  Further, these
schedules are always permutations of the rule-instance set, i.e., each
rule-instance appears in the schedule exactly once.  The dynamic
semantics has no such restriction.

As described earlier in Sec.~\ref{sec_schedules}, this semantics takes
no position on schedules.  A schedule is a parameter to this
semantics, and may therefore be produced by an oracle, even one that
dynamically selects the next rule at the last moment.  There is no
such thing in this semantics as a ``wrong'' schedule; different
schedules merely admit more or less concurrency (because they'll
encounter less or more conflicts, respectively).

The \emph{bsc} compiler chooses a schedule by static analysis of the
rules in a program.  In other words, it is statically predicting
something about dynamic execution (rule and method conditions, rule
conflicts and hardware conflicts).  This semantics has exact
information, but the compiler can only have estimated or approximate
information.

Thus, there are rule firings that may be permitted in this dynamic
semantics that will be prohibited in code compiled by \emph{bsc},
because \emph{bsc} can only use conservative approximations (e.g.,
\emph{bsc} may see a potential conflict where the dynamic semantics
knows there is none).

Thus, the semantics described here is more liberal than what is
produced today by \emph{bsc}, and covers possible future improvements
in \emph{bsc} schedule analysis and compilation.

% ================================================================

\vspace{1cm}

\hdivider

\bibliographystyle{abbrv}
\bibliography{BSV_Semantics}

\end{document}
